\documentclass[11pt,a4paper]{article}
\usepackage[margin=23mm]{geometry}
\usepackage{amsmath,amssymb,booktabs,array}
\input{glyphtounicode}
\pdfgentounicode=1
\usepackage[hidelinks]{hyperref}
\setlength{\parindent}{0pt}
\setlength{\parskip}{6pt}
\begin{document}
\begin{center}
{\Large H1: a homogeneous upstream candidate}\\[4pt]
13 October 2026 --- analytical qualification, not an adopted traffic contract
\end{center}

\section*{1. Traffic meaning and assumption boundary}
The ending lane $E$ seeks early transfer into adjacent continuing lane $A$;
continuing lane $B$ does not receive lateral traffic. A finite uniform
footprint initially occupies $[0,\ell]$ inside a fixed preparation area
$[0,L]$, with $0<\ell\leq L$. At the decision instant it has densities
$e_0,a_0,b_0$ and mass $M=\ell(e_0+a_0+b_0)>0$. Initial symmetry means
$a_0=b_0$, not $a(t)=b(t)$ after merging. No new arrivals are included.
All quantities and the initial footprint are fixed across action comparisons.

One constant common advice $v\in[v_-,v_+]$, $v_->0$, is realized as the
preparation speed on the admitted free-flow branch. This is a declared
analytical response hypothesis, not established compliance. Lane changes
are represented by a uniform same-coordinate lateral source. The model does
not reproduce the diagonal next-cell movement of the native CTM algorithm.

The physical interaction scale $h>0$ is fixed across actions. It converts a
throughput-equivalent opportunity into a lateral density-change rate. Its
meaning is repeated physical opportunity per unit progress, rather than an
arbitrary numerical tick. It is an explicit adaptation requiring provenance
or robustness qualification; no source-derived or fitted value is adopted.
The early work weight $\gamma\geq1$ and late work weight $\alpha\geq1$ are
distinct. Work is normalized so that an ordinary mass unit uses one work unit.

\begin{center}
\begin{tabular}{lll}
\toprule
Quantity & Meaning & Units\\
\midrule
$L,\ell,h,z$ & geometry, footprint, opportunity scale, progress & m\\
$v,w$ & common free speed, backward-wave magnitude & m/s\\
$k_j,e,a,b$ & jam or lane density & veh-eq/m\\
$q_E,q_A,q_B$ & ordinary longitudinal mass flux & veh-eq/s\\
$C,S$ & receiving-work rate (ordinary-work normalization) & work/s\\
$f$ & lateral mass-flux equivalent at a representative opportunity & veh-eq/s\\
$m=f/h$ & local lateral density source & veh-eq/(m s)\\
$\gamma,\alpha$ & receiving work per lateral/late mass unit & dimensionless\\
\bottomrule
\end{tabular}
\end{center}

The boundary must accept approaching traffic for the no-wait transport
qualification below. A waiting terminal is considered as a reopen condition,
not silently added to the equations.

\section*{2. State to receiving, own demand, and completed merging}
Use one declared triangular diagram for the diagnostic family:
\[
 Q_v(k)=\min\{vk,w(k_j-k)\},\quad
 k_c(v)=\frac{wk_j}{v+w},\quad C(v)=v k_c(v),\qquad w,k_j>0.
\]
For density $a$ the adjacent receiving rate is
\[
 S(v,a)=\min\{C(v),w(k_j-a)\}.
\]
Restrict the initial densities over the whole action range by
\[
 0\leq e_0,a_0,b_0\leq k_c(v_+).
 \tag{A}
\]
The candidate will keep $a\leq k_c(v)$ and $e\leq e_0$, so ordinary traffic
remains on the linear free-flow branch. Own longitudinal flow and the declared
rate analogue of eager donor availability are
\[
 q_A=va,\quad q_E=ve,\quad q_B=vb_0,\qquad D_E=ve.
\]
Within this branch $S=C$ and $q_A\leq S$. The local R1-inspired rate allocation
is
\[
 f(v,e,a)=\min\left\{D_E,\frac{S(v,a)-q_A}{\gamma}\right\},\qquad
 m=\frac{f}{h}.
 \tag{1}
\]
Consequently
\[
 0\leq f\leq ve,\qquad va+\gamma f\leq S(v,a),\qquad
 \dot e=-m,\quad \dot a=m,\quad \dot b=0.
 \tag{2}
\]
Actual transferred mass is $m\,dx\,dt$; it is not $\gamma m\,dx\,dt$.
No clipping rule is needed inside the admitted domain because $S-q_A\geq0$.

\textbf{Adaptation boundary.} Carey's R1 supports own-lane-first residual
receiving and the gap-weighted work inequality. Its native request is a
cell stock in a native interval. Equations (1)--(2) additionally choose a
progress-based rate offer, same-coordinate transfer and a physical scale $h$.
They are not a verbatim continuous-time transcription of Algorithm A1. Here
$q_E$ describes road-frame longitudinal motion; $f$ is the opportunity-equivalent
bound used in a separate lateral source. It is not a second longitudinal
outflow to add to $q_E$ at the same boundary. A native diagonal donor sending
budget cannot be imported unchanged into this altered representation.

Likewise, $C$ at successive road cross sections is not one pooled capacity
to sum and spend remotely. Uniformity makes the local budget identical at
every occupied location; the choice $f/h$ specifies lateral intensity locally.
This is precisely the additional homogeneous interaction assumption to qualify.

\section*{3. Why the homogeneous member is invariant}
For an analytical transport reference, use the proposed balance equations
within the occupied region:
\[
 \partial_t e+v\partial_x e=-m(e,a;v),\qquad
 \partial_t a+v\partial_x a=m(e,a;v),\qquad
 \partial_t b+v\partial_x b=0.
 \tag{3}
\]
Outside the common footprint all three densities are zero and the source is
zero. Equal initial values at all occupied positions produce equal solutions
of the same ordinary differential equations along their characteristics.
Thus the occupied body stays uniform while its edges translate at $v$.
Equation (3) is a reference proof, not a position-resolved simulator or an
assertion about congested waves in this task.

Put $r=e_0+a_0=e+a$. In the free-flow member (2) becomes
\[
 \dot e=-\frac{v}{h}\min\left\{e,\frac{k_c(v)-r+e}{\gamma}\right\},
 \qquad a=r-e.
 \tag{4}
\]
At $e=0$ the rate is zero. The donor-bound $m\leq ve/h$ preserves
nonnegative ending density. At $a=k_c(v)$ the spare receiving vanishes, so
the adjacent density cannot cross that boundary. Therefore
\[
 0\leq e(t)\leq e_0,\qquad
 a_0\leq a(t)\leq\min\{r,k_c(v)\},\qquad b(t)=b_0.
 \tag{5}
\]
Condition (A) and this invariant ensure the selected common speed is feasible
throughout the candidate upstream motion. They do not prove that every
traffic state under a VSL actually lies in this domain.

\section*{4. The opportunity clock and exact branch references}
With physical progress $z=vt$ and $\delta(v)=k_c(v)-r$, (4) gives
\[
 \frac{de}{dz}=-\frac1h\min\left\{e,\frac{\delta(v)+e}{\gamma}\right\}.
 \tag{6}
\]
The explicit factor $v$ cancels. If normalized receiving $C(v)/v$ is
unchanged between actions, this candidate has the same preparation per unit
distance. Additional elapsed time alone supplies no benefit.

For a donor-limited branch starting at $(z_s,e_s)$,
\[
 e(z)=e_s\exp\left[-\frac{z-z_s}{h}\right].
 \tag{7}
\]
For a residual-service branch,
\[
 e(z)=(e_s+\delta)\exp\left[-\frac{z-z_s}{\gamma h}\right]-\delta.
 \tag{8}
\]
Use these formulas only while the corresponding minimum branch remains active.
For $\gamma>1$ and $\delta>0$, the switch threshold is
$e=\delta/(\gamma-1)$. If $\delta<0$, the donor cannot vanish while the
target stays free flow; its limiting residual is $-\delta=r-k_c(v)$.
These are exact analytical branch references, not numerical rollouts.

In the triangular diagnostic
\[
 k_c'(v)=-\frac{wk_j}{(v+w)^2}<0.
\]
For a service branch with fixed initial branch data,
\[
 \frac{\partial e}{\partial\delta}
 =\exp\left[-\frac{z-z_s}{\gamma h}\right]-1\leq0.
\]
Slower speed can therefore increase preparation via changed normalized
density headroom. This conditional preparation sensitivity is not the
derivative of complete time with a waiting terminal. Switch events and any
new boundary need their own derivative qualification later.

\section*{5. Compact physical arrivals and whole-cohort conservation}
The front reaches $L$ at $t_f=(L-\ell)/v$, and the rear at $t_r=L/v$.
For $0\leq t<t_r$, the length still in preparation is
\[
 s(t)=\min\{\ell,L-vt\};\qquad s(t)=0\text{ for }t\geq t_r.
\]
The no-wait arrival rates at the lane end are
\[
 A_E'(t)=ve(t),\quad A_A'(t)=va(t),\quad A_B'(t)=vb_0,
 \qquad t_f<t<t_r,
 \tag{9}
\]
and zero outside that interval. These are arrivals to terminal service,
not guaranteed complete departures when service binds. Ordinary total
arrival rate is $v\rho_0$, with $\rho_0=e_0+a_0+b_0$.

Before the front exits, remaining total mass is $s(t)\rho_0=M$.
During exit, it is $s(t)\rho_0$ and decreases by the total arrival rate.
Including all terminal arrivals, conservation is
\[
 s(t)\rho_0+ A_E(t)+A_A(t)+A_B(t)=M.
 \tag{10}
\]
For the ending lane in particular, early transferred mass is
$\int_0^t s(u)m(u)\,du$, and
\[
 s(t)e(t)+A_E(t)+\int_0^t s(u)m(u)\,du=\ell e_0.
 \tag{11}
\]
Hence the moving-footprint formulation does not lose mass when its occupied
length shrinks. It also does not create an arbitrary current-stock discharge
tail. Every surviving upstream element reaches the boundary according to its
initial physical position and common travel clock.

\section*{6. A joint nonblocking terminal limit and its scientific conclusion}
Let $R_A^{\rm out}(v,t)$ and $R_B^{\rm out}(v,t)$ be the declared downstream
receiving-work rates. For a nonblocking terminal it is sufficient that
\[
 va(t)+\alpha ve(t)\leq R_A^{\rm out}(v,t),\qquad
 vb_0\leq R_B^{\rm out}(v,t)
 \tag{12}
\]
throughout the arrival window for all admitted actions. These are domain
conditions, not externally supplied future receiving curves for prediction.
As one existence member, keep declared constant finite downstream receiving
rates $Q_A^{\rm out},Q_B^{\rm out}>0$ at an unconstrained exit. This retains
the fixed-resource idea of S13 without adopting its simultaneous batch or
proportional allocation. Conservative sufficient initial bounds are
\[
 v_+\alpha r\leq Q_A^{\rm out},\qquad
 v_+b_0\leq Q_B^{\rm out}.
 \tag{13}
\]
Indeed $a+\alpha e=r+(\alpha-1)e\leq\alpha r$, and
$v\leq v_+$. Sufficiently sparse positive initial traffic can satisfy (A)
and (13) together. This member has no terminal queue and no feedback
that contradicts the upstream homogeneous clock. Allocation priority is
immaterial when every demand is served; no new binding-terminal priority is
selected by this limit.

Every initial position $x\in[0,\ell]$ completes at $(L-x)/v$. Its mean
complete time and derivative are therefore
\[
 \overline T(v)=\frac{L-\ell/2}{v},\qquad
 \frac{d\overline T}{dv}=-\frac{L-\ell/2}{v^2}<0.
 \tag{14}
\]
The preferred speed is the largest admissible one, irrespective of hidden
merge composition. The result is a limiting/negative information-dependent
control finding. It does not transfer S13's positive information/action
certificate to the new member.

If (12) fails, terminal waiting is consequential. A queue on the ending lane
or adjacent lane can change incoming flow and earlier receiving. Neither
(3) nor (9) then automatically remains valid. A detached point queue with
unlimited storage and no upstream influence would be a new approximation or
new geometry assumption, not a source-supported default. This limitation does
not prove that all compact homogeneous holding models are impossible.

For any eventual admitted terminal, retain the complete-population objective
\[
 F(v)=\int_0^{T_{\rm clear}} N_{\rm unexited}(t;v)\,dt,
 \qquad \overline T=F/M.
 \tag{15}
\]
It counts upstream travel, any preparation waiting, terminal holding and all
traffic until final exit. Staggered arrivals do not justify importing the
half-clearance mean of a simultaneous proportional batch.

\section*{7. Information sufficiency and present disposition}
With fixed known footprint and initial symmetry, aggregate mass $M$ plus
initial ending mass $E_0$ gives
\[
 e_0=E_0/\ell,\qquad a_0=b_0=(M-E_0)/(2\ell).
\]
Equation (4) then determines $e(t)$, and $a(t)=e_0+a_0-e(t)$ reconstructs
the adjacent receiver. Thus one initial ending statistic is sufficient for
the proposed upstream member. A strong aggregate-plus-ending-statistic rival
can implement the identical equations and histories. Independent initial
continuing-lane asymmetry would change that information contract; it is not
admitted here.

\textbf{Disposition.} Homogeneous upstream closure is mathematically possible
under explicit adaptations. A complete no-blocking member exists but loses
the desired information-dependent speed contrast. A waiting-terminal member,
applied MP6 domain, fresh reachable observer witness, derivative protocol and
independent usefulness evidence remain open. The review stops at the original
B2 system-interface gate; no simulator or experiment is authorized.

\section*{Source and adaptation record}
\begin{itemize}
\item Carey, Balijepalli and Watling (2013), \emph{Extending the Cell
Transmission Model to Multiple Lanes and Lane-Changing}: diagram and local
allocation ingredients, not this joined homogeneous task.
\href{https://eprints.whiterose.ac.uk/id/eprint/77315/7/Carey\%20et\%20al\%202013_with_coversheet.pdf}{Author manuscript}, Sections 2, 4 and 5.
\item Jin (2015), \emph{Continuous formulations and analytical properties of
the link transmission model}: initial/boundary-profile requirements, not
one-stock closure with interior lateral sources.
\href{https://arxiv.org/pdf/1405.7080}{Primary preprint}, Section 2.2.
\item Jin (2021), \emph{A link queue model of network traffic flow}:
compact model-form rival with distinct internal dynamics.
\href{https://arxiv.org/pdf/1209.2361}{Primary preprint}, introduction and conclusion.
\end{itemize}
All additional profile, rate-scale, uniform-source, speed-response and terminal
assumptions are declared above. Conditional proof is not empirical calibration.
The accepted S13 contract and its original scalar certificate are unchanged.
\end{document}
